\documentclass[10pt]{article}

% Specify the margins and text size.
\setlength{\textwidth}{6.4in}
\setlength{\textheight}{9.5in}
\setlength{\oddsidemargin}{0pt}
\setlength{\evensidemargin}{0pt}
\setlength{\topmargin}{0pt}
\setlength{\hoffset}{.05in}
\setlength{\voffset}{-1in}

\setlength{\parskip}{5pt}
\setlength{\parindent}{0pt}

% Load some fonts and symbol packages
\usepackage{latexsym}
\usepackage{pifont}       % contains 'star' symbol for counterinsurgency handbook title
\usepackage{yfonts} 
\usepackage{amsmath}
\usepackage{amsfonts}

\usepackage{graphicx}     % actually, this is loaded by pstricks
\usepackage[T1]{fontenc}
\usepackage{ifthen}
\usepackage{pstricks,pst-grad,pst-text,pst-node,multido,pst-plot,calc,pst-3dplot}
%\usepackage[all]{xy}
%\usepackage{animate}

% The hyperref package inserts links into the pdf.
\definecolor{MyLinkColor}{rgb}{.1,.2,1}
\definecolor{MyCiteColor}{rgb}{.1,1,.2}
\definecolor{MyURLColor}{rgb}{.4,.4,.4}
\usepackage[backref=true,pagebackref=false,hyperindex,colorlinks=true,
  linkcolor=MyLinkColor,urlcolor=MyURLColor]{hyperref}


% The tweaklist package is something I found on the web.  It provides a simple interface
% for making changes to spacing used in the itemize and enumerate environments.  Comment
% this out if you don't care to use tweaklist.
\usepackage{tweaklist}
\renewcommand{\itemhook}{\setlength{\parskip}{2pt}\setlength{\parsep}%
{1pt}\setlength{\topsep}{0pt}\setlength{\itemsep}{0pt}}

\newcommand{\U}{\underline{\hspace{5pt}}}

\usepackage{listings}
\newcommand{\Z}{\hphantom{0}}
\newcommand{\HH}{\hspace{20pt} \hphantom{a)} }

\begin{document}
\pagestyle{empty}
\lstset{language=R, showspaces=false, showstringspaces=false}

\href{http://www.shepherd.edu}{\includegraphics[height=1.75cm]{logo-high-res.eps}}
\vspace{-1.79cm}

{\small
\begin{tabular}{cl}
& Math 314\\
& Statistics\\
\hspace{5.28in} & 9 November 2011
\end{tabular}
}
\setlength{\baselineskip}{1.05\baselineskip}

\begin{center}
\textbf{\large  Quiz 7}
\end{center}

1. With a perfectly balanced roulette wheel, in the long run, 
red numbers should turn up 18 times in 38.
To test its wheel, one casino records the results of 3,800 plays, finding 1,890 red
numbers.  Is that too many reds?  Or chance variation?

\hspace{20pt} a) Formulate the null and alternative hypotheses as 
statements about a box model.
\vspace{1.5in}

\hspace{20pt} b) The null says that the percentage of reds in the box is
\underline{\hspace{30pt}}.  The alternative says that the\vspace{-4pt}

\HH percentage of reds in the box is 
\underline{\hspace{30pt}}.
\bigskip

\hspace{20pt} c) Compute $z$ and $P$.
\vspace{3.5in}

\hspace{20pt} d) Were there too many reds?  Justify your answer.

\vfill
\eject
{\ }

2. One kind of plant has only blue flowers and white flowers.  According to a 
genetic model, the offsprings of a certain cross have a 75\% chance
to be blue-flowering.  Two hundred seeds of such a cross are raised, and
142 turn out to be blue-flowering.  Are the data consistent with the model?  
Answer yes or no, and explain briefly.
\vfill
\eject

\end{document}

